\section{What the Laboratory Already Knows}

Watch what a measurement actually involves, and count the things known in the course of it that no instrument reports.

Begin with the arithmetic. A physicist averages five readings; the average requires that five is what the readings number, that addition behaves tomorrow as it behaved today, and that division does not depend on the operator's mood. No experiment established these truths, and none could overturn them: an apparent laboratory refutation of arithmetic would be treated---rightly---as a miscount, a malfunction, or a joke. The same holds for logic. The whole evidential force of an experiment runs through inference: \emph{if} the theory is true, \emph{then} the meter reads thus; the meter does not read thus; therefore the theory is in trouble. That pattern---the logician's \emph{modus tollens}---is not a finding of any science. It is what makes findings possible. A mind that had to verify logic empirically could never get an empirical verification started, since the bearing of evidence on hypothesis \emph{is} a logical relation.\endnote{An empiricist of Quine's stripe replies that logic and mathematics are not known apart from experience but are simply the most protected members of a single web of belief that faces experience as a whole, revisable in principle like everything else (W.~V. Quine, ``Two Dogmas of Empiricism,'' \emph{Philosophical Review} 60 [1951]: 20--43; cited from the standard bibliographic record, not collated for this essay). The reply is serious, but notice what it concedes and what it costs: it concedes that nothing is verifiable sentence-by-sentence, which was the criterion's whole mechanism; and any deliberate \emph{revision} of the web must itself be conducted by inference---some norm of consistency decides what ``facing experience'' requires---so the authority of logic reappears inside the very act said to revise it. Either way, the instrumentalist's sharp test is gone.}

Continue with the reach of the conclusions. Science does not want yesterday's meter readings; it wants laws---claims about all bodies, all charges, every case, observed and unobserved, past and future. The inference from \emph{tested here} to \emph{holds everywhere} is called induction, and its warrant is famously not itself an observation. Hume, the patron of the flames, was also the man who showed that no experience can certify that unexperienced cases resemble experienced ones, since any argument from past success of induction to its future success is induction again.\endnote{Hume, \emph{Enquiry}, sect.~IV. The point here is not Humean skepticism---this essay takes induction to be rationally warranted---but the location of the warrant: whatever justifies trusting nature's uniformity, it is not delivered by an instrument, since instruments report actual cases only and the uniformity claim outruns every actual case.} The working scientist is right to trust induction. But the trust is a piece of rational confidence that outruns every possible instrument reading, held by every experimentalist as a condition of experiment.

Continue with the observations themselves, for even they are not the criterion's innocent atoms. What counts as a datum---which flicker is signal, which is noise, which run is discarded for a known fault---is decided under a theory, by a trained judgment the instrument cannot supply. The point is not a modern subtlety; the genealogy found it in Comte's own founding chapter: ``facts cannot be observed without the guidance of some theory\ldots\ for the most part we could not even perceive them.''\endnote{Comte, in Martineau's condensation, vol.~1, p.~4, as identified and checked in the genealogy above. Comte used the point to explain why primitive man needed theology; it applies with equal force to the fantasy of a tribunal of bare observation before which all theory is arraigned. The tribunal's evidence is theory-soaked from the first exhibit---which does not make observation circular or untrustworthy, but does make ``observation alone'' a description of nothing.} The father of positivism, explaining history, stated the premise that dissolves positivism's picture of a self-certifying observational ground. Observation disciplines theory magnificently---but only from inside a practice that already knows more than observation delivers.

Continue with the sociology of evidence. Almost nothing any scientist knows was verified by that scientist. She knows the mass of the electron from tables, the reliability of the tables from the practice of the community, the results of the crucial experiments from reports written by the dead. Her own data from last year she knows by memory and by trusting records she has no independent means to check against the past itself. Testimony, memory, the existence of other minds competent to testify---this is the load-bearing structure of scientific knowledge, and none of it registers on any gauge. A consistent instrumentalist, permitted to believe only what he could bring before observation, would be reduced to the contents of his present visual field; he could not even believe in the previous page of his own lab notebook, since the past that wrote it is exactly the kind of thing no present observation can contain.

Continue, finally, with the norms. The entire enterprise runs on the conviction that data ought to be reported honestly, that a beautiful theory must yield to an ugly fact, that fabricating results is not merely imprudent but \emph{wrong}. These are judgments about obligation, and obligation casts no shadow on a detector. On Ayer's official account they are not judgments at all but expressions of feeling, so that ``falsifying data is wrong'' says no more than a groan says. No scientific institution could survive a week on that reading of its own rules, and none has ever tried. The instrumentalist's laboratory is held together by moral knowledge his theory classifies as noise; his position is, in the exact sense, parasitic---alive only because the knowledge it denies keeps feeding it.

And at the center of all of it sits the thing the instruments cannot be: the one who reads them. A needle against a scale is not information; it is brass and ink until a mind takes it \emph{as} standing for a quantity, under a theory, within a calibration, toward a conclusion. Every measurement terminates in an act of understanding and judgment that is not itself a measurement. This is not mysticism about consciousness; it is the plainest description of what happens when anyone reads a dial. The instrument extends the senses magnificently. It does not extend one inch beyond the senses toward the judgment that makes sense of them.

Now assemble the inventory. Inside the most disciplined empirical practice humanity has built, we find knowledge of mathematics and logic, which no observation establishes; rational warrant for induction, which outruns all observation; knowledge by testimony and memory of a past and of other minds that no observation contains; moral knowledge of the norms of inquiry; and the knower's own understanding, presupposed by every reading and reported by none. None of this is exotic. All of it is knowledge, if anything is---the instrumentalist himself stakes his practice on every item. The thesis that only the instrument-testable can be known is therefore not a modest generalization from science; it is a description that fits no science, no scientist, and no single hour of scientific work that has ever been performed.

What follows is not that anything goes. The right conclusion is more precise, and it is the hinge of this essay: knowledge is \emph{proportioned to its objects}. Quantities are known by measuring because quantity is the sort of thing measurement fits. Theorems are known by proof, particulars of the past by testimony and record, obligations by practical reason, minds by their expression in speech and act---and in each case the method answers to the kind of thing being known, and is answerable to it, testable in its own register, correctable by its own failures. This plurality is the opposite of relativism: it is rigor, distributed. The demand that every object submit to the method proper to quantity is not rigor but a category mistake enforced with a confident face.

One question, then, can no longer be dismissed on procedural grounds. Sensible, measurable, contingent things exist; the sciences know them with growing precision. Whether such things point beyond themselves---whether the existence of the measurable has a ground that is not itself one more measurable thing---is a question about reality, askable in earnest, condemned by no coherent criterion. It must be judged the way every serious question is judged: by the arguments. To those arguments, and to the tradition that spent two thousand years refining them, the essay now turns.
